\chapter{Safety Requirements}

FR-08 -- Modular Mechanical Interface The robot shall include a mechanical mounting area allowing future modules to be installed.

4. Safety Requirements

This section defines the safety mechanisms, constraints, and operational behaviors required to ensure safe deployment of the AMR-X mobile robot in indoor warehouse environments. The system must prioritize human safety, prevent collisions, ensure controlled motion, and guarantee safe failure handling under all operating conditions

a) Emergency Stop

\begin{itemize}
    \item The robot shall include a physical emergency stop button accessible on the robot
body.

    \item The system shall provide an emergency stop function via the operator dashboard.
    \item Activation of the emergency stop shall result in immediate motor power cutoff.
    \item The emergency stop shall override all ongoing operations and motion commands.
\end{itemize}

\begin{itemize}
    \item The system shall automatically trigger an emergency stop in case of critical sensor
failure (e.g., LiDAR failure, IMU failure, or loss of odometry).

\end{itemize}

b) Obstacle Avoidance and Collision Prevention

\begin{itemize}
    \item The robot shall use LiDAR-based perception for real-time obstacle detection.
    \item The system shall perform autonomous collision avoidance during navigation.
    \item The robot shall reduce speed in congested or high-risk environments.
    \item The robot shall stop immediately if an unrecoverable collision risk is detected.
    \item The system shall maintain a safety buffer distance defined in the navigation costmap.
c) Manual Override and Control

    \item The operator shall be able to pause autonomous navigation at any time.
    \item The operator shall be able to resume or cancel ongoing missions.
    \item The system shall support a manual control mode (implementation method TBD:
joystick or keyboard interface).

    \item All manual override actions shall be logged with timestamp and operator identity for
traceability.

\end{itemize}

d) Speed Limitation and Control

\begin{itemize}
    \item The robot shall enforce a configurable maximum speed per operational sector.
    \item The system shall automatically reduce speed in areas with detected human presence
or obstacles nearby.

    \item Speed profiles shall be adaptable based on environmental mapping and risk level.
    \item Specific numeric speed thresholds per sector remain TBD and will be defined during
system design phase.

\end{itemize}

e) Low Battery Safety Behavior

\begin{itemize}
    \item The robot shall continuously monitor battery levels during operation.
    \item A low battery warning shall be triggered at 20\% capacity.
    \item At 10\% battery level, the robot shall initiate an autonomous return-to-docking
behavior.

    \item If battery level becomes critical during a mission, the system shall pause or safely
terminate the mission. (we can add an integrating charging station like the one developed by HAI Robotics or a dedicated zone can be defined where robots with low battery levels are required to navigate for recharging or safe standby.)

\end{itemize}

f)    Mission Failure Handling

\begin{itemize}
    \item In case of navigation failure, the robot shall stop and notify the operator with an
error code.

    \item In case of communication loss, the robot shall enter a safe state and attempt return-
to-dock behavior if possible.

    \item In case of sensor failure, the system shall either:
\end{itemize}

degrade functionality safely, or perform an emergency stop depending on severity.

\begin{itemize}
    \item All failures shall be logged with diagnostic information for debugging and
traceability.

\end{itemize}

g) Restricted Zones

\begin{itemize}
    \item The system shall support geofencing capabilities to define restricted or forbidden
areas.

    \item Navigation planning shall avoid restricted zones automatically.
    \item Missions attempting to pass through restricted areas shall be rejected by the system.
    \item Manual override of restricted zones may be allowed but must be logged and
authorized.

    \item Environmental constraints shall be enforced using Nav2 costmap-based planning (or
equivalent system).

\end{itemize}

h) Human Safety

\begin{itemize}
    \item The robot shall continuously monitor its environment for human presence detection.
    \item The robot shall reduce speed when operating in shared human-robot environments.
    \item The system shall include visual indicators to show operational state (idle, moving,
error, autonomous mode).

    \item The robot shall provide audio or visual warnings during autonomous movement.
\end{itemize}

i)    Security Requirements

\begin{itemize}
    \item All dashboard access shall require secure authentication (login system).
    \item The system shall enforce role-based access control (RBAC) (roles to be defined in
future design phase).

    \item The robot shall not expose any unauthenticated control endpoints.
    \item All operator actions shall be recorded in an audit log with timestamps and user
identification.

    \item Communication between system components shall use encrypted channels.
\end{itemize}

